% Full-page figure: high-dimensional volume concentration.
% Two stacked plots:
%   (a) Fraction of n-ball volume within outer shell of width
%       epsilon, as a function of n, for epsilon = 0.01, 0.05, 0.1.
%   (b) Ratio of inscribed-ball volume to enclosing-cube volume,
%       V_n(1)/2^n, on a log scale, n = 1..50.
\begin{figure}[p]
\centering
\begin{tikzpicture}[
  every axis/.style={
    font=\small,
    axis lines=left,
    width=12cm,
    height=6.5cm,
    grid=both,
    grid style={engblack!12, very thin},
    major grid style={engblack!22},
    enlargelimits=false,
    tick style={engblack},
    label style={font=\small},
  },
]

% ---- Panel (a): outer-shell fraction vs dimension n ----
\begin{axis}[
  name=shellplot,
  xlabel={dimension $n$},
  ylabel={fraction in outer shell $1 - (1-\varepsilon)^{n}$},
  xmin=0, xmax=500,
  ymin=0, ymax=1.02,
  legend pos=south east,
  legend cell align=left,
  legend style={font=\footnotesize, draw=engblack!30},
  title={(a) Volume of $n$-ball inside outer shell of width $\varepsilon R$},
]
  % epsilon = 0.01
  \addplot[engblue, very thick, domain=1:500, samples=200] {1 - (0.99)^x};
  \addlegendentry{$\varepsilon = 0.01$}
  % epsilon = 0.05
  \addplot[engred, very thick, dashed, domain=1:500, samples=200] {1 - (0.95)^x};
  \addlegendentry{$\varepsilon = 0.05$}
  % epsilon = 0.10
  \addplot[engblack, very thick, dotted, domain=1:500, samples=200] {1 - (0.90)^x};
  \addlegendentry{$\varepsilon = 0.10$}
  % 50% reference line
  \addplot[engblack!40, thin, domain=0:500] {0.5};
\end{axis}

% ---- Panel (b): inscribed-ball / cube ratio, log scale ----
\begin{axis}[
  name=ratioplot,
  at={(shellplot.below south west)},
  anchor=above north west,
  yshift=-1.4cm,
  xlabel={dimension $n$},
  ylabel={$V_{n}(1) / 2^{n}$ (log scale)},
  xmin=1, xmax=50,
  ymode=log,
  ymin=1e-30, ymax=1.5,
  title={(b) Ratio of inscribed unit ball to enclosing cube $[-1, 1]^{n}$},
]
  % V_n(1)/2^n = pi^(n/2) / (2^n Gamma(n/2 + 1))
  % Use the recursive log form is awkward in pgfplots; precompute as
  % a coordinates table. Values from V_n(1) = pi^(n/2)/Gamma(n/2+1).
  \addplot[engblue, very thick, mark=*, mark size=1.2pt] coordinates {
    (1, 1.0)
    (2, 0.7854)
    (3, 0.5236)
    (4, 0.3084)
    (5, 0.1645)
    (6, 0.0807)
    (7, 0.0369)
    (8, 0.0159)
    (9, 0.00641)
    (10, 0.00249)
    (12, 3.26e-4)
    (15, 1.43e-5)
    (20, 2.46e-8)
    (25, 1.84e-11)
    (30, 6.65e-15)
    (40, 3.28e-22)
    (50, 5.10e-30)
  };
\end{axis}

\end{tikzpicture}
\caption[High-dimensional volume concentration in the ball and the
cube.]{(a) The fraction of the $n$-dimensional ball that sits in the
outer shell of relative thickness $\varepsilon$ is $1 - (1 -
\varepsilon)^{n}$. For $\varepsilon = 0.01$ (solid), the fraction
crosses $50\%$ near $n \approx 70$ and exceeds $99\%$ by $n
\approx 500$. (b) The volume of the unit ball inscribed in the cube
$[-1, 1]^{n}$ collapses to a negligible fraction of the cube as $n$
grows: the ratio is $\pi/4 \approx 0.79$ at $n = 2$ and drops below
$10^{-8}$ by $n = 20$. Together the two plots are the working
illustration of the failure section. Random sampling in
high-dimensional geometry should not assume that a "typical" sample
lies near the centre of the domain or even near the inscribed ball.}
\label{fig:vol02:ch11:volume-concentration}
\end{figure}
